%! TEX root = ../bcd-19-20.tex

\chapter{Șiruri de funcții}

\section{Convergență punctuală și convergentă uniformă}

Fie $ (f_n) $ un șir de funcții, adică fiecare $ f_n : D \seq \RR \to \RR $
este o funcție reală. Dacă în cazul șirurilor de numere reale, limita este un
număr real, în cazul șirurilor de funcții, limita este o funcție.

Însă noțiunea de convergență este mai fină în cazul șirurilor de funcții.

Astfel, avem două tipuri de convergență:
\begin{definition}\label{def:conv-punct}\index{șiruri!funcții}\index{convergență!punctuală}
    Fie $ (f_n) $ un șir de funcții reale ca mai sus.

    Spunem că șirul $ (f_n) $ \emph{converge punctual (simplu)} la funcția $ f $
    dacă are loc:
    \[
        \lim_{n \to \infty} f_n(x) = f(x), \quad \forall x \in D.
    \]
    Vom nota aceasta pe scurt prin $ f_n \xrar{pc} f $ sau $ f_n \xrar{s} f $,
    iar funcția $ f $ se va numi \emph{limita punctuală} a șirului $ (f_n) $.
\end{definition}

De asemenea, un alt tip de convergență de interes este \emph{convergența uniformă}.
Formal, ea se definește folosind criterii de caracterizare cu $ \varepsilon $, dar în
exerciții vom folosi cel mai des teorema de mai jos.
\begin{theorem}\label{thm:conv-unif}\index{convergență!uniformă}
    În condițiile și cu notațiile de mai sus, spunem că șirul $ (f_n) $ converge
    uniform la funcția $ f $, notat $ f_n \xrar{uc} f $ sau $ f_n \xrar{u} f $
    dacă:
    \[
        \lim_{n \to \infty} \sup_{x \in D} |f_n(x) - f(x)| = 0.
    \]
\end{theorem}

Evident, funcția $ f $ folosită în convergența uniformă va fi limita punctuală
a șirului, deci este clar că proprietatea de convergență uniformă o implică
pe cea de convergență punctuală.

Proprietatea reciprocă este falsă, după cum se vede din exemplul simplu
de mai jos.

Fie $ f_n : [0, 1] \to \RR, f_n(x) = x^n, \forall n \geq 1 $.

Atunci se poate vedea imediat că:
\[
    \lim_{n \to \infty} f_n(x) =%
    \begin{cases}
        0, & 0 \leq x < 1 \\
        1, & x = 1
    \end{cases}.
\]
Rezultă că $ f_n \xrar{s} f $, unde $ f $ este funcția definită pe ramuri
mai sus.

Un calcul simplu arată însă că:
\[
    \lim_{n \to \infty} \sup_{x \in [0, 1]} |f_n(x) - f(x)| = 1 \neq 0,
\]
deci șirul nu este uniform convergent.

%%%%%%%%%%%%%%%%%%%%%%%%%%%%%%%%%%%%%%%%%%%%%%%%%%%%%%%%%%%%%%%%%%%%%%
\section{Transferul proprietăților analitice}

În continuare, ne punem problema următoare: presupunem că fiecare termen
al șirului de funcții $ (f_n) $ are o proprietate analitică de un anumit
tip (privitoare la continuitate, derivabilitate, integrabilitate etc.). 
Întrebarea este care dintre aceste proprietăți se transferă și funcției
limită (punctuală sau uniformă) și în ce condiții.

În tot ceea ce urmează, vom păstra notațiile și ipoteza de mai sus,
adică vom lucra cu șirul de funcții $ (f_n) $, cu fiecare
$ f_n : D \seq \RR \to \RR $.

\begin{theorem}[Transfer de continuitate]\label{thm:transf-cont}
    Dacă fiecare $ f_n $ este funcție continuă, pentru orice $ n $,
    iar șirul $ (f_n) $ converge uniform la funcția $ f $, atunci
    funcția $ f $ este continuă.
\end{theorem}

Proprietatea de mai sus poate fi folosită ca o condiție suficientă
în felul următor: dacă limita simplă a unui șir de funcții nu este
continuă, atunci șirul nu poate fi uniform convergent. Aceasta deoarece
funcția care este limita simplă a șirului este singurul candidat pentru
convergența uniformă.

\begin{theorem}[Integrare termen cu termen]\label{thm:transf-int}
    Dacă $ f_n \xrar{u} f $, atunci are loc proprietatea de integrare
    termen cu termen, adică:
    \[
        \lim_{n \to \infty} \int_a^b f_n(x) dx = \int_a^b f(x) dx, \quad%
        \forall [a, b] \seq D.
    \]
\end{theorem}

Cu alte cuvinte, în ipoteza convergenței uniforme, limita și integrala pot comuta.

\begin{theorem}[Derivare termen cu termen]\label{thm:transf-der}
    Presupunem că funcțiile $ f_n $ sînt derivabile, pentru orice $ n \in \NN $.
    Dacă $ f_n \xrar{s} f $ și dacă există o funcție $ g : D \seq \RR \to \RR $
    astfel încît $ f'_n \xrar{u} g $, atunci $ f $ este derivabilă și $ f' = g $.
\end{theorem}

%%%%%%%%%%%%%%%%%%%%%%%%%%%%%%%%%%%%%%%%%%%%%%%%%%%%%%%%%%%%%%%%%%%%%%
\section{Exerciții}

1. Să se studieze convergența punctuală și uniformă a șirurilor de funcții:
\begin{enumerate}[(a)]
\item $ f_n : (0, 1) \to \RR, f_n(x) = \dfrac{1}{nx + 1}, n \geq 0 $;
\item $ f_n : [0, 1] \to \RR, f_n(x) = x^n - x^{2n}, n \geq 0 $;
\item $ f_n : \RR \to \RR, f_n(x) = \sqrt{x^2 + \dfrac{1}{n^2}}, n > 0 $;
\item $ f_n : [-1,1] \to \RR, f_n(x) = \dfrac{x}{1 + nx^2} $;
\item $ f_n : (-1, 1) \to \RR, f_n(x) = \dfrac{1 - x^n}{1 - x} $;
\item $ f_n : [0, 1] \to \RR, f_n(x) = \dfrac{2nx}{1 + n^2 x^2} $;
\item $ f_n : \RR_+ \to \RR, f_n(x) = \dfrac{x + n}{x + n + 1} $;
\item $ f_n : \RR_+ \to \RR, f_n(x) = \dfrac{x}{1 + nx^2} $.
\end{enumerate}

\vspace{1cm}

2. Să se arate că șirul de funcții dat de:
\[
  f_n : \RR \to \RR, \quad f_n(x) = \frac{1}{n} \arctan x^n
\]
converge uniform pe $ \RR $, dar:
\[
  \Big( \lim_{n \to \infty} f_n(x) \Big)'_{x = 1} \neq \lim_{n \to \infty} f'_n(1).
\]

Rezultatele diferă deoarece șirul derivatelor nu converge uniform pe $ \RR $.

\vspace{1cm}

3. Să se arate că șirul de funcții dat de:
\[
  f_n : [0, 1] \to \RR, \quad f_n(x) = nxe^{-nx^2}
\]
este convergent, dar:
\[
  \lim_{n \to \infty} \int_0^1 f_n(x) dx \neq \int_0^1 \lim_{n \to \infty} f_n(x) dx.
\]
Rezultatul se explică prin faptul că șirul nu este uniform convergent.
De exemplu, pentru $ x_n = \dfrac{1}{n} $, avem $ f_n(x_n) \to 1 $, dar,
în general, $ f_n(x) \to 0 $.

\vspace{1cm}

